%% Route bsi-counter, round 16.  Paste-ready; amsthm style of frontier.tex.
%% Labels prefixed bsc:.  Everything below is PROVED unless marked GAP.
%% Full file: /tmp/claude-1000/-data-ssg-proof/d1fe2115-9b72-4784-bb94-87421ac1106c/scratchpad/bsi-counter/bsc.tex

\subsection{Best-response restart: what its rounds are made of}

The rule $R_{\mathrm{BR}}$ of \Cref{prop:bsi-br} is stated through a veto, and
the veto is opaque.  It has a transparent form.  For Max strategies
$\sigma,\sigma'$ write $\mathrm{Dis}(\sigma,\sigma'):=\{v\in\Vmax:\sigma(v)\ne\sigma'(v)\}$,
and for a Min strategy $\tau$ let
$\mathrm{BR}(\tau):=\{\sigma':\val_{\sigma',\tau}=\val^{\tau}\}$ be the set of
Max best responses to $\tau$.

\begin{lemma}[$R_{\mathrm{BR}}$ moves towards the double best response]\label{bsc:double-br}
Let $G$ be stopping, $\tau$ a Min strategy and $U:=\val^{\tau}$.
\begin{enumerate}
\item[(a)] $\sigma'\in\mathrm{BR}(\tau)$ if and only if $\sigma'$ is
  \emph{$U$-greedy}: $U(\sigma'(v))=\max_iU(v^{(i)})$ for every $v\in\Vmax$.
\item[(b)] For every Max strategy $\sigma$ and every $v\in\Vmax$,
\[
  U(\bar\sigma(v))>U(\sigma(v))\iff \sigma'(v)\ne\sigma(v)\ \text{for all }
  \sigma'\in\mathrm{BR}(\tau),\qquad
  U(\bar\sigma(v))\ge U(\sigma(v))\iff \sigma'(v)=\bar\sigma(v)\ \text{for some }
  \sigma'\in\mathrm{BR}(\tau).
\]
\end{enumerate}
Consequently, if $\tau$ is a Min best response to $\sigma$, the two variants of
\Cref{prop:bsi-br} switch
\[
  C^{>}_{\max}=S_\sigma\cap\!\!\bigcap_{\sigma'\in\mathrm{BR}(\tau)}\!\!\mathrm{Dis}(\sigma,\sigma'),
  \qquad
  C^{\ge}_{\max}=S_\sigma\cap\!\!\bigcup_{\sigma'\in\mathrm{BR}(\tau)}\!\!\mathrm{Dis}(\sigma,\sigma'),
\]
and when the Max best response $\sigma^{\ast}$ to $\tau$ is unique both are
$S_\sigma\cap\mathrm{Dis}(\sigma,\sigma^{\ast})$: one round of $R_{\mathrm{BR}}$
switches exactly those improving Max vertices at which $\sigma$ disagrees with
a best response to a best response to $\sigma$.
\end{lemma}

\begin{proof}
(a) $U=\val^{\tau}$ is the unique fixed point of $T^{\tau}$
(\Cref{thm:contraction}, $G$ stopping), so $U(v)=\max_iU(v^{(i)})$ at
$v\in\Vmax$, $U(v)=U(\tau(v))$ at $v\in\Vmin$ and $U$ is the mean of its
successors at average vertices.  If $\sigma'$ is $U$-greedy then $U$ satisfies
all the equations of the profile $(\sigma',\tau)$, and by uniqueness of that
fixed point $U=\val_{\sigma',\tau}$.  Conversely if $\val_{\sigma',\tau}=U$
then $U(v)=U(\sigma'(v))$ at every $v\in\Vmax$, so $\sigma'(v)$ attains
$\max_iU(v^{(i)})$.  (b) is immediate from
$U(v)=\max\{U(v^{(0)}),U(v^{(1)})\}$ and (a).
\end{proof}

\begin{remark}\label{bsc:double-br-rem}
\Cref{bsc:double-br} makes \Cref{prop:bsi-br} a statement about best-response
dynamics: $R_{\mathrm{BR}}$ is the map
$\sigma\mapsto\sigma[S_\sigma\cap\mathrm{Dis}(\sigma,\sigma^{\ast\ast})]$ with
$\sigma^{\ast\ast}$ a best response to a best response to $\sigma$, the
intersection with $S_\sigma$ being what makes it a switching rule
(\Cref{def:rule}) rather than an unguarded double best response, which does
not converge.  It also explains \Cref{thm:readout}: on $\mathrm{RD}(n)$ the
double best response to $\sigma_t$ disagrees with $\sigma_t$ at
$v_{t+1}$ only.
\end{remark}

\begin{lemma}[A constant option is taken at most twice]\label{bsc:const}
Let $G$ be stopping and let $\sigma_0,\sigma_1,\dots$ be produced by any
switching rule (\Cref{def:rule}), so $\sigma_{t+1}=\sigma_t[S_t]$ with
$\emptyset\ne S_t\subseteq S_{\sigma_t}$.  Let $v\in\Vmax$ have successors
$p,q$ with $\val_{\sigma}(q)=c$ for every Max strategy $\sigma$ --- for
instance $q$ a vertex from which no controlled vertex is reachable.  Then $v$
is switched at most twice along the whole run: at most once from $p$ to $q$,
at most once from $q$ to $p$, and in that order.  Dually for a Min vertex with
a constant successor along the Min track of \Cref{def:bsi}.
\end{lemma}

\begin{proof}
By \Cref{lem:switch} the vector $\val_{\sigma_t}$ is pointwise nondecreasing in
$t$, so $\varphi(t):=\val_{\sigma_t}(p)-c$ is nondecreasing.  If
$\sigma_t(v)=p$ then $v\in S_{\sigma_t}$ iff $\varphi(t)<0$; if
$\sigma_t(v)=q$ then $v\in S_{\sigma_t}$ iff $\varphi(t)>0$.  Since $\varphi$
is nondecreasing, $\{t:\varphi(t)<0\}$ is an initial segment and
$\{t:\varphi(t)>0\}$ a final segment of the run, and they are disjoint.  Every
$p\to q$ switch occurs in the first set and every $q\to p$ switch in the
second, so all $p\to q$ switches precede all $q\to p$ switches; two $p\to q$
switches would need a $q\to p$ switch between them, which is impossible.
\end{proof}

\begin{corollary}\label{bsc:const-cor}
If every $v\in\Vmax$ has a constant successor then every switching rule halts
within $2|\Vmax|$ rounds on $G$; so no game of that shape is a hard instance
for any switching rule, $R_{\mathrm{BR}}$ and all-switches included.
\end{corollary}

\begin{lemma}[Min's preference at a clamp reverses once]\label{bsc:minrev}
Let $\sigma_0,\sigma_1,\dots$ be a run of a switching rule and let
$u\in\Vmin$ have successors $p,q$ with $\val_\sigma(q)=c$ constant.  A Min
best response $\tau_t$ to $\sigma_t$ has $\tau_t(u)=p$ when
$\val_{\sigma_t}(p)<c$ and $\tau_t(u)=q$ when $\val_{\sigma_t}(p)>c$.  Hence
if ties are broken by a fixed preference at $u$, the action $\tau_t(u)$ is a
monotone step function of $t$ and changes at most once along the run.
\end{lemma}

\begin{proof}
$\val_{\sigma,\tau}=\val_\sigma$ forces $\tau(u)$ to attain
$\min_i\val_\sigma(u^{(i)})$ (dual of \Cref{bsc:double-br}(a)), and
$t\mapsto\val_{\sigma_t}(p)-c$ is nondecreasing by \Cref{lem:switch}.
\end{proof}

\begin{theorem}[Blocks, and the reordering bound]\label{bsc:reorder}
Consider a run of $R_{\mathrm{BR}}$, in either variant, on a stopping $G$ with
$|\Vmax|=m$, producing $\sigma_0,\dots,\sigma_R$ with guides
$U_t:=\val^{\tau_t}$.  Put
\[
  d_t:=\bigl|\{v\in\Vmax:\ U_{t-1}(\sigma_t(v))=\textstyle\max_iU_{t-1}(v^{(i)})
   \ \text{and}\ U_t(\sigma_t(v))<\max_iU_t(v^{(i)})\}\bigr|\qquad(1\le t\le R-1),
\]
the number of Max vertices at which the guide's preference \emph{reverses away
from the current action} at round $t$.  Then
\[
  R\ \le\ m+\sum_{t=1}^{R-1}d_t .
\]
In particular $d_t=0$ whenever $U_t=U_{t-1}$, so a maximal block of
consecutive rounds sharing one guide has at most $m$ rounds and
$R\le m\bigl(1+\#\{t:\tau_t\ne\tau_{t-1}\}\bigr)$.  For \Cref{def:bsi}, in
either variant, the same argument gives
\[
  R\ \le\ m+B+\sum_{t}d_t\qquad\text{and dually}\qquad R\ \le\ k+A+\sum_{t}d'_t,
\]
with $A,B$ as in \Cref{thm:bsi-tracks}; since $d_t\le m$ and $d_t=0$ at a
round with $C_{\min}=\emptyset$, this refines \Cref{thm:bsi-tracks}'s
$R\le m+(m+1)B$.
\end{theorem}

\begin{proof}
Put $\Phi(\sigma,U):=|\{v\in\Vmax: U(\sigma(v))=\max_iU(v^{(i)})\}|\in[0,m]$.
Every $v\in C_{\max}$ has $U_t(\bar\sigma_t(v))\ge U_t(\sigma_t(v))$, so
$\bar\sigma_t(v)$ is $U_t$-greedy and $\sigma_{t+1}(v)$ is $U_t$-greedy;
unswitched vertices keep their status.  Hence
$\Phi(\sigma_{t+1},U_t)=\Phi(\sigma_t,U_t)+a_t$, where $a_t:=|C^{>}_{\max}|$
counts the switched vertices that were \emph{not} $U_t$-greedy.  Next,
$\Phi(\sigma_{t+1},U_{t+1})\ge\Phi(\sigma_{t+1},U_t)-d_{t+1}$ by the
definition of $d_{t+1}$.  Summing over $t$ and using $0\le\Phi\le m$,
\[
  \sum_{t=0}^{R-1}a_t\ \le\ m+\sum_{t=1}^{R-1}d_t .
\]
For $R_{\mathrm{BR}}$, \Cref{prop:bsi-br} in its strict reading gives
$C^{>}_{\max}\ne\emptyset$ at every round with $S_\sigma\ne\emptyset$, so
$a_t\ge1$ and $R\le\sum a_t$.  For \Cref{def:bsi}, $a_t\ge1$ can fail; but by
\Cref{lem:bsi-pairloc} in its strict reading every productive round has
$C^{>}_{\max}\ne\emptyset$ or $C^{<}_{\min}\ne\emptyset$, and the rounds of
the second kind have $C_{\min}\ne\emptyset$ and so number at most $B$.
\end{proof}

\begin{corollary}[Clamped games are easy for $R_{\mathrm{BR}}$]\label{bsc:clamp}
Let $G$ be stopping with $|\Vmax|=m$, $|\Vmin|=k$, and suppose every Min
vertex has a constant successor.  Then $R_{\mathrm{BR}}$, run with any fixed
tie-breaking preference at each Min vertex, halts within $m(k+1)$ rounds from
every start.  Dually, if every Max vertex has a constant successor, every
switching rule halts within $2m$ rounds (\Cref{bsc:const-cor}).
\end{corollary}

\begin{proof}
By \Cref{bsc:minrev} each Min vertex changes its chosen action at most once,
so $\tau_t$ changes at most $k$ times and the run splits into at most $k+1$
blocks; each has at most $m$ rounds by \Cref{bsc:reorder}.
\end{proof}

\begin{remark}\label{bsc:clamp-rem}
\Cref{bsc:clamp} is why every family this document carries is fast for
$R_{\mathrm{BR}}$: $\mathrm{RD}(n)$ has all its $\Theta_j$ constant
($m=k=n$, bound $n(n+1)$, truth $n$), $L_n$, $H_m$, $\mathrm{CC}(L,m)$,
$\mathrm{WD}(e,j,m)$ and $\mathrm{SD}(K)$ have $k=0$, and
\Cref{prop:bsi-oneplayer} is the case $k=0$ of the corollary.  Together with
\Cref{bsc:const} it says what a superpolynomial family must contain:
\emph{both} a Max vertex whose two options are non-constant and a Min vertex
whose two options are non-constant, the Min one reversing superpolynomially
often.  This subsumes the round-$15$ auditor's refutation of the $R_{\mathrm{BR}}$
route's proposed counter (a Max vertex over a constant gadget) and gives its
Min-side twin.
\end{remark}

\subsection{The ceiling of $R_{\mathrm{BR}}$ at two Max vertices}

Write $h_{\mathrm{BR}}(m)$ for the greatest number of productive rounds of
$R_{\mathrm{BR}}$ over all stopping SSGs with $|\Vmax|=m$, all starts and both
variants --- the analogue for this rule of the four ceilings of
\Cref{rem:four-ceilings}.  Since $\val_{\sigma}$ strictly increases at every
productive round (\Cref{lem:switch}), no Max strategy recurs, so $R\le2^{m}-1$
for $R_{\mathrm{BR}}$ and, for \Cref{def:bsi}, $R\le(2^{m}-1)+(2^{k}-1)$.
Both bounds are attained at $m=2$, $k=1$.

\begin{proposition}[$W_2$: three rounds with two Max vertices]\label{bsc:w2}
Let $W_2$ be the stopping SSG whose harmonic normal form
(\Cref{lem:normalform}) has $\Vmax=\{v_0,v_1\}$, $\Vmin=\{u\}$ and rows
\[
\begin{array}{r@{\ }l@{\qquad}l}
 v_0\ (\max): & \tfrac{13}{16}\,v_1 & \tfrac1{32}v_1+\tfrac{23}{32}u+\tfrac5{32}t_1\\
 v_1\ (\max): & \tfrac{29}{32}\,u   & \tfrac14\,t_1\\
 u\ (\min):   & \tfrac{15}{16}\,v_0 & \tfrac1{32}\,t_1
\end{array}
\]
(the missing mass leaks to $t_0$; every row is dyadic of denominator $32$, so
\Cref{lem:dyadic-row} realises $W_2$ as a stopping SSG on $34$ vertices: $2$
Max, $1$ Min, $29$ average and the two sinks).  Then
$w^{\ast}=(\tfrac{13}{64},\tfrac14,\tfrac1{32})$ and:
\begin{enumerate}
\item[(i)] from $\sigma_0=(0,0)$, $R_{\mathrm{BR}}$ takes exactly three
 productive rounds, $(0,0)\to(1,0)\to(1,1)\to(0,1)$, switching $v_0$, then
 $v_1$, then $v_0$ again; the run is the same in the veto and the strict
 variant and for \emph{every} choice of best response, and no comparison made
 along it is a tie;
\item[(ii)] all-switches takes at most two rounds from every start, and
 $h^{\ast}(2)=2$ (\Cref{prop:auso-census}), so $R_{\mathrm{BR}}$ beats the
 all-switches ceiling already at $m=2$;
\item[(iii)] $R=3=2^{m}-1$ is the absolute maximum for $m=2$, so
 $h_{\mathrm{BR}}(2)=3$ against $h^{\ast}(2)=2$, $h^{\ast}_{HK}(2)=2$;
\item[(iv)] \Cref{def:bsi} from $(\sigma_0,\tau_0)=((0,0),0)$ takes exactly
 four rounds in both variants, with switch sets
 $(\{v_0\},\emptyset),(\emptyset,\{u\}),(\{v_1\},\emptyset),(\{v_0\},\emptyset)$,
 and $4=(2^{m}-1)+(2^{k}-1)$ is the absolute maximum for $(m,k)=(2,1)$.  In
 particular a BSI run can be longer than $|C|$, which the two families of
 \Cref{rem:readout} are not.
\end{enumerate}
The two guides involved are
$\val^{\tau=0}=(\tfrac{2560}{4909},\tfrac{2175}{4909},\tfrac{2400}{4909})$ with
Max best response $(1,0)$, and $\val^{\tau=1}=w^{\ast}$ with Max best response
$(0,1)$; the Min best response is $\tau=0$ at $\sigma_0$ and $\tau=1$ at every
later strategy, and the value vectors along the run are
$\val_{\sigma_0}=(0,0,0)$,
$\val_{\sigma_1}=(\tfrac{5885}{32768},\tfrac{29}{1024},\tfrac1{32})$,
$\val_{\sigma_2}=(\tfrac{191}{1024},\tfrac14,\tfrac1{32})$ and
$\val_{\sigma_3}=w^{\ast}$.
\end{proposition}

\begin{proof}
The rows all leak, so $W_2$ is stopping by \Cref{lem:trapchar}.  Everything
else is a finite exact computation, carried out twice: in the normal form,
with $\val_\sigma$ the componentwise minimum over both Min strategies and
$\val^{\tau}$ the componentwise maximum over all four Max strategies; and on
the realised $34$-vertex SSG with the independent core of the harness, whose
$w^{\ast}$ and whose $R_{\mathrm{BR}}$ run agree.  For (iii) and the second
half of (iv): $\val_\sigma$ strictly increases at each round with
$C_{\max}\ne\emptyset$ and $\val^{\tau}$ strictly decreases at each round with
$C_{\min}\ne\emptyset$ (\Cref{lem:switch} in $G$ and in $\bar G$), so no
$\sigma$ and no $\tau$ recurs and there are at most $2^{m}-1$ rounds of the
first kind and $2^{k}-1$ of the second, every productive round being of one
kind or the other.
\end{proof}

\begin{remark}\label{bsc:w2-rem}
The mechanism is worth naming, because \Cref{bsc:clamp} says every family must
eventually break it.  The Min vertex $u$ is a \emph{clamp},
$\val_\sigma(u)=\min(\tfrac{15}{16}\val_\sigma(v_0),\tfrac1{32})$, and it
engages once, as soon as Max's own improvement pushes $\val_\sigma(v_0)$ past
$\tfrac1{30}$.  Before it engages, the guide $\val^{\tau}$ is computed in a
counterfactual world in which Min does \emph{not} clamp, so $u$ is worth
$\tfrac{2400}{4909}$ there instead of $\tfrac1{32}$; that inflated $u$ makes
$v_1$ prefer reading $u$ to its constant $\tfrac14$, which is the veto that
serialises round $0$, and makes $v_0$ prefer the option that reads $u$, which
is the switch that round $0$ performs.  After the clamp engages the guide is
$w^{\ast}$ and the last two rounds are ordinary greedy steps.  So the veto's
selectivity comes from the \emph{gap between the guide's world and the true
one at a Min vertex that has not yet been forced}, and by \Cref{bsc:minrev} a
clamp affords exactly one such gap.
\end{remark}

\subsection{How many guides a one-switch-per-round run needs}

\begin{proposition}[The piercing bound]\label{bsc:pierce}
Let a run of $R_{\mathrm{BR}}$ on a stopping $G$ with $|\Vmax|=m$ switch
exactly one Max vertex per round and visit $\sigma_0,\dots,\sigma_R$, and let
$f_t$ be the vertex switched at round $t$.  Then
\begin{enumerate}
\item[(a)] $S_{\sigma_t}=\{c:\val_{\sigma_t[c]}>\val_{\sigma_t}\}$, so if
 $\sigma_t[c]$ is itself visited, $c\in S_{\sigma_t}$ iff it is visited later;
\item[(b)] the double best response $\sigma^{\ast}_t$ of \Cref{bsc:double-br}
 lies in the subcube
 $C_t:=\{\sigma^{\ast}:\sigma^{\ast}(c)=\sigma_t(c)\ \forall c\in
 S_{\sigma_t}\setminus\{f_t\},\ \sigma^{\ast}(f_t)\ne\sigma_t(f_t)\}$, and
 $\sigma_R\in C_R:=\{\sigma_R\}$;
\item[(c)] hence the number of distinct Min best responses $\tau_t$ occurring
 along the run is at least the piercing number $p$ of the family
 $\{C_t\}$, and $k\ge\log_2p$.
\end{enumerate}
For the run that visits all $2^{m}$ strategies in the reflected Gray order ---
the order of the least-index run on the ladder $L_m$ of \Cref{def:ladder},
verified here for $m\le4$ --- the piercing numbers are
\[
 m=2,3,4,5,6,7:\qquad p=2,\,4,\,6,\,\le9,\,\le14,\,\le22 ,
\]
the first three exact by exhaustive search and the last three by a randomised
greedy cover; and at $m=3$ \emph{every} Hamiltonian run needs at least three
guides (exhaustively over all $144$ Hamiltonian paths of the $3$-cube; $48$ of
them need $3$ and $96$ need $4$).  So a $2^{m}-1$-round $R_{\mathrm{BR}}$ run
needs $k=\Omega(m)$ Min vertices but no more --- the piercing count is
exponentially far below $2^{m}$ and is not an obstruction.
\end{proposition}

\begin{proof}
(a) is \Cref{lem:switch} together with its converse: if $c\notin S_\sigma$
then $T_{\sigma[c]}\val_\sigma\le\val_\sigma$ pointwise, so
$\val_{\sigma[c]}\le\val_\sigma$ by \Cref{cor:comparison}.  Since
$\val_{\sigma_t}$ strictly increases along the run, a visited neighbour is
improving iff it is visited later.  (b) is \Cref{bsc:double-br}: the round
switches $S_{\sigma_t}\cap\mathrm{Dis}(\sigma_t,\sigma^{\ast}_t)=\{f_t\}$,
which says exactly that $\sigma^{\ast}_t$ agrees with $\sigma_t$ on
$S_{\sigma_t}\setminus\{f_t\}$ and differs at $f_t$; at the halt $\sigma_R$ is
optimal and its best response $\tau_R$ is optimal too, so
$\sigma_R\in\mathrm{BR}(\tau_R)$.  (c) A guide $\tau$ determines
$\sigma^{\ast}$, so distinct rounds with the same guide have the same
$\sigma^{\ast}$, and the set of $\sigma^{\ast}$'s used pierces $\{C_t\}$.  The
numbers are computations.
\end{proof}

\subsection{What is left, stated exactly}

\begin{remark}[The measurements]\label{bsc:measured}
All in exact rational arithmetic, both variants, and (where stated) over every
choice of best response.
\begin{enumerate}
\item[(i)] The inequalities of \Cref{bsc:reorder} (for $R_{\mathrm{BR}}$ and
 for \Cref{def:bsi}), the block bound, \Cref{bsc:const} and \Cref{bsc:clamp}
 were checked on $600$ freshly generated random normal forms with
 $m\in\{2,3,4\}$, $k\in\{1,2,3\}$ and denominators $8,16$, from every start
 and (for BSI) every start pair, in both variants, and on $600$ random
 \emph{clamped} games: no violation.  The bound $R\le m+\sum d_t$ is attained.
 The longest $R_{\mathrm{BR}}$ run met by random sampling was $4$ and the
 longest BSI run $5$; on the random clamped games $R\le m$ always, so
 \Cref{bsc:clamp}'s $m(k+1)$ is nowhere near met by sampling.
\item[(ii)] $\mathrm{RD}(n)$ of \Cref{def:readout-cascade} has $R=n$ with $n$
 blocks of one round each, $\sum d_t=n-1$ and every Min preference reversing
 once; $L_n$ has one block; $W_2$ has $R=3$ in two blocks of $1$ and $2$
 rounds and $\sum d_t=2$; $M_3$ has $R=4$ in two blocks of $1$ and $3$ rounds,
 so a block that is a full cascade of $m$ single switches under one frozen
 guide --- the first half of what \textsc{Gap 1} needs --- is realised
 already at $m=3$, and there too the single Min preference reverses once.
\item[(iii)] Records found by exact hill-climbing on dyadic normal forms
 (denominator $16$), maximising the round count over all starts: $R=3$ at
 $m=2,k=1$ (\Cref{bsc:w2}, the maximum possible), $R=4$ at $m=3,k=1$ on the
 game $M_3$ recorded in the route's files ($52$ vertices when realised;
 $h^{\ast}(3)=4$, so $R_{\mathrm{BR}}$ has not been separated from
 all-switches at $m=3$ here); at $m=4$ nothing above $4$ was found in the
 time available, and the document's own $G^{\sharp}$ already gives $5$.  A
 five-round run at $m=3$ was looked for and not found: by \Cref{bsc:pierce}
 eighteen of the thirty single-switch runs of length five in the $3$-cube can
 be guided by two targets and hence by one Min vertex, and each was prescribed
 to a float annealer (several hundred restarts, $k=1$ and $k=2$), as was a
 hand-shaped fourteen-parameter family with two clamps; the exact hill-climb
 and an exhaustive neighbourhood search at denominators $16$ and $64$ leave
 $M_3$ a strict local maximum.  Failure to find is not evidence of absence,
 and this document's own record is that every hard instance in it had to be
 engineered rather than searched for.
\end{enumerate}
\end{remark}

\begin{remark}[What a superpolynomial family must contain]\label{bsc:programme}
Collecting \Cref{bsc:const,bsc:minrev,bsc:reorder,bsc:clamp}:
\begin{enumerate}
\item[(a)] every Max vertex switched more than twice has two non-constant
 options, and every Min vertex whose best response reverses more than once has
 two non-constant options: the round-$15$ auditor's refutation and its exact
 Min-side twin;
\item[(b)] $R\le m+\sum_td_t$, so the guide must \emph{reorder} the options at
 Max vertices $\Omega(R)$ times in total --- changing $\tau$ is not enough,
 the change must reverse a preference;
\item[(c)] in the clamped class --- every Min vertex with a constant successor
 --- $R\le m(k+1)$, which is $O(N^{2})$ and no more.  So the clamped class
 cannot carry a superpolynomial family, and the most a clamp construction can
 give is quadratic.
\end{enumerate}
Two statements would be real progress, and neither is proved here.

\medskip\noindent\textsc{Gap 1 (quadratic, clamped).} \emph{There is a family
of stopping SSGs on $O(n\log n)$ vertices with $|\Vmax|=|\Vmin|=\Theta(n)$, in
which every Min vertex has a constant successor, on which $R_{\mathrm{BR}}$
takes $\Theta(n^{2})$ rounds from a stated start.}  By (c) this is the
ceiling of the clamped class; it would be the first superlinear lower bound
for the rule, $\mathrm{RD}(n)$ giving only $\Theta(N)$.  The shape is forced
by \Cref{bsc:reorder}: $k$ clamps engaging in succession give $k+1$ guides,
and each block must be a cascade of $m$ single switches, so the construction
needs a Max chain that sweeps forwards under one guide and backwards under the
next.

\medskip\noindent\textsc{Gap 2 (superpolynomial).} \emph{There is a stopping
SSG with a Min vertex $u$, both of whose options are non-constant, along whose
$R_{\mathrm{BR}}$ run the sign of
$\val_{\sigma_t}(u^{(0)})-\val_{\sigma_t}(u^{(1)})$ changes superpolynomially
often.}  By \Cref{bsc:reorder} this is necessary; by \Cref{prop:leapfrog} the
analogous reversal at a \emph{Max} vertex is available at $O(1)$ vertices per
reversal, but there the reversals are driven by fresh drivers that switch once
each, and what is needed here is a leapfrog whose drivers are the run itself.
\end{remark}